\chapter{Background}
\label{ch:background}
% 2026-11-25 새 장: 지도교수 요청(배경·관련연구 분량을 본문의 약 25%로). 옛 2장의 ATT&CK 절
% (sec:related-attack)을 이리로 옮겨 넓혔다. 새 인용의 원문 사본은 sources/papers/<key>.pdf,
% 초고와 인용 표는 sources/expand/{sec,mlbg}.md (저장소에 올리지 않음).
% 2026-11-26 지도교수 지시(본문 40~50쪽, 배경·관련연구는 비율 유지로 12쪽 안팎): 7쪽을 4쪽으로 줄였다.
%   줄이기 전 판은 git 기록(커밋 3303d09)에 있다. 남긴 문장의 원문 대조 주석은 그대로 두었다.

This chapter collects the general concepts that the later chapters
assume: the evidence, Windows audit records of process creation and
termination, and the vocabularies of adversary behavior
(Sections~\ref{sec:bg-audit} and~\ref{sec:related-attack}), and on the
machine-learning side the scoring of probabilistic forecasts, seed
randomness, ensembles and distillation, and distribution shift
(Sections~\ref{sec:bg-scoring} to~\ref{sec:bg-shift}). The definitions
of this thesis itself are given in Chapters~\ref{ch:task}
to~\ref{ch:evaluation}.

\section{Windows host auditing}
\label{sec:bg-audit}

Windows records security-relevant activity in the Security event log
when the corresponding audit policy is enabled. Audit Process Creation
generates event 4688 each time a new process
starts~\cite{msauditprocesscreation,ms4688}, and Audit Process
Termination generates event 4689 each time a process
exits~\cite{ms4689}.
% 원문: Microsoft Learn 4688 "Subcategory: Audit Process Creation", "This event generates every time a new
%   process starts."; 4689 "Subcategory: Audit Process Termination", "every time a process has exited". 확인됨.
An event 4688 record names the account under which the creating
process ran, with a logon identifier, the identifier and full path of
the new process (\emph{New Process Name}), the identifier and path of
its creator, the token elevation type, the integrity label, and the
\emph{Process Command Line}~\cite{ms4688}. The command line is filled
only when the policy \emph{Include command line in process creation
events} is also enabled, and is then written in plain
text~\cite{ms4688,mscmdline}. An event 4689 record names the account,
the identifier and path of the exited process, and an
application-defined exit status~\cite{ms4689}.
% 원문: Microsoft Learn 4688 필드 설명(Creator Subject Logon ID, New Process Name, Creator Process ID,
%   Token Elevation Type, Mandatory Label); "By default Process Command Line field is empty."; "Command line
%   process auditing" "logged in plain text"; 4689 "Exit Status ... unique for every application". 확인됨.
Matching the Creator Process ID of a 4688 record with the New Process
ID of another, and the Process ID of a 4689 record with a New Process
ID on the same computer, yields a tree of
processes~\cite{ms4688,ms4689}. The link is made by the analyst or the
collector; the records are independent log entries.

Records leave the host through a collection path. Windows Event
Forwarding sends selected events to a collector in batches, every 15
minutes in its built-in \emph{Normal} option and every 30 seconds with
\emph{Minimize Latency}, and buffers them in the local log while the
collector is unreachable; events overwritten before then are lost
without notice~\cite{mswef}.
% 원문: Microsoft Learn "Use Windows Event Forwarding to help with intrusion detection"(2025-08-18) FAQ
%   "batch timeout of 15 minutes", "batch timeout of 30 seconds", "local event log is the buffer",
%   "there's no notification sent to the WEF collector that events are lost". 확인됨.
The time at which a collector holds a record is therefore set by
batching, buffering, and connectivity, not only by the time the event
occurred; Chapter~\ref{ch:task} builds its input definition on this.
In research, host audit records have mainly served to explain or
detect behavior after it has occurred
(Section~\ref{sec:related-provenance}); this thesis reads them to
forecast an operation that has not yet started.

\section{Adversary behavior models}
\label{sec:related-attack}

MITRE ATT\&CK organizes observed adversary behavior into tactics, the
objectives an activity serves, and techniques and sub-techniques, the
ways an objective is pursued~\cite{mitreintro,mitrewindows}. It began
as a categorization of post-compromise behavior against Windows
systems to improve detection, and its authors state that it is not an
exhaustive enumeration of attack vectors; procedures are the specific
implementations of a technique, and one procedure can span several
techniques~\cite{strom2018mitre}.
% 원문: MITRE ATT&CK "Get Started"(2026-11-24 열람) Tactics/Techniques/Sub-techniques; Strom et al. (2018,
%   rev. 2020) §1 p.1 "categorize post-compromise ... against Microsoft Windows systems to improve detection",
%   "not an exhaustive enumeration of attack vectors"; §3.4.1 p.9 "a procedure can span multiple techniques". 확인됨.
Our 18 prediction labels (Appendix~\ref{app:actions}) are finer than
ATT\&CK techniques and name operations at the level of the controlled
procedures. A system information enumeration and a registry query are
both discovery activity, related to System Information Discovery and
Query Registry respectively; merging them would erase the distinction
the classifier is asked to make. The label is the operation designated
by the procedure and verified against the actual process start, not an
ATT\&CK identifier.

The intrusion kill chain divides an intrusion into seven phases, from
reconnaissance to actions on objectives, so that defenders can reason
about where a mitigation breaks the chain~\cite{hutchins2011killchain}.
% 원문: Hutchins et al. §3.2 p.4 일곱 단계 정의, 1절 "just one mitigation disrupts the chain". 확인됨.
The Diamond Model links the events against one victim into activity
threads, phase-ordered graphs of attack paths that actually
occurred~\cite{caltagirone2013diamond}, and Attack Flow links actions,
conditions, and relationships into a connected flow of
techniques~\cite{mitreattackflow}.
% 원문: Caltagirone et al. §8 p.30 "An activity thread is a directed phase-ordered graph", p.34 "activity threads
%   model knowledge of actual attack paths"; CTID Attack Flow "actions, conditions, and relationships as a
%   connected flow". 확인됨.
These models describe relations within an attack sequence, are written
by an analyst from activity already known, and use units coarser than
one host operation: a phase, a tactic, a technique, or a Diamond event.
At a single host decision several of our operations fall under the
same phase or tactic, and the ATT\&CK authors likewise note that
lifecycle models do not convey individual actions or how one relates
to another~\cite{strom2018mitre}.
% 원문: Strom et al. §4.1.3 p.22 "not effective at conveying what individual actions adversaries make, how one
%   action relates to another". 확인됨.
This thesis keeps the sequence view, but its forecaster receives only
the prefix that arrived before the next action and must commit to a
distribution over that action before it is observed.

\section{Probabilistic next-action prediction and its scoring}
\label{sec:bg-scoring}

A probabilistic classifier over $K$ classes maps an input $x$ to a
distribution $p(\cdot\mid x)$; here the classes are the $K=18$ actions
and the input is the record received up to a prediction boundary
(Chapter~\ref{ch:task}). A neural classifier computes logits
$z(x)\in\mathbb{R}^K$ and applies the softmax at temperature $T>0$,
\begin{equation}
  p_T(a\mid x)=\frac{\exp\bigl(z_a(x)/T\bigr)}
                    {\sum_{b=1}^{K}\exp\bigl(z_b(x)/T\bigr)} ,
  \label{eq:bg-softmax}
\end{equation}
where $T=1$ is the ordinary softmax and $T$ does not change the
ranking of the classes.

A scoring rule is proper if a forecaster whose belief is $q$ obtains
the best expected score by issuing $q$, and strictly proper if $q$ is
the only such forecast~\cite{gneiting2007strictly}. Written as losses,
the two classical strictly proper rules for a categorical outcome $y$
are the logarithmic and the Brier loss,
\begin{equation}
  \ell_{\log}(p,y)=-\log p(y\mid x),
  \qquad
  \ell_{\mathrm{B}}(p,y)=\sum_{a=1}^{K}\bigl(p(a\mid x)-\indicator[a=y]\bigr)^2 .
  \label{eq:bg-scores}
\end{equation}
% 원문: Gneiting & Raftery (JASA 2007) 초록 p.359 proper/strictly proper; §3.1 p.363 Examples 1–3 strictly proper,
%   Example 3 Bregman distance = KL, "criticized for its unboundedness"; Example 4 zero–one "proper but not
%   strictly proper". 확인됨. 손실 방향(부호 뒤집기)은 우리가 쓴 표기다.
For the logarithmic loss the excess expected loss of $p$ over the true
distribution $q$ is $\mathrm{KL}(q\,\|\,p)$~\cite{gneiting2007strictly}.
The loss is unbounded, so a few confident errors can dominate a mean.
Negative log-likelihood (NLL) is the mean of $\ell_{\log}$ over
evaluation rows; ours weights each execution equally
(Equation~\ref{eq:nll}). Top-1 accuracy, the zero--one score, is proper
but not strictly proper~\cite{gneiting2007strictly}: forecasts with the
same mode score the same however they spread the rest of their
probability, so two models can agree in accuracy and differ in NLL.

Calibration asks whether stated confidence matches observed frequency.
Guo et al.\ found modern neural networks poorly calibrated and fitted a
single $T$ in Equation~\ref{eq:bg-softmax} to validation NLL after
training, which leaves the predicted class and the accuracy
unchanged~\cite{guo2017calibration}.
% 원문: Guo et al. (ICML 2017) 초록 p.1 "poorly calibrated"; §4.2 p.6 "T is optimized with respect to NLL on the
%   validation set", "does not affect the model's accuracy". 확인됨.
NLL rewards calibration and sharpness together. This thesis scores
models by NLL and accuracy (Section~\ref{sec:eval-metrics}); post hoc
temperature scaling enters only as a comparison arm
(Section~\ref{sec:gen-methods}), and the temperature $\tau$ of
LogitNorm in our training loss is a different quantity
(Section~\ref{sec:method-baseline}).

\section{Training randomness and underspecification}
\label{sec:bg-seeds}

Training a neural network is a randomized procedure. Its sources of
nondeterminism include parameter initialization, the shuffling of data
into minibatches, augmentation, stochastic regularization such as
dropout, and low-level library operations, which a seed cannot control
but only switch on or
off~\cite{summers2021nondeterminism,bouthillier2021accounting}.
% 원문: Summers & Dinneen (ICML 2021) §3 p.2 목록, "it is not possible to 'seed' this"; Bouthillier et al.
%   (MLSys 2021) §2.2 p.3 다섯 변동원. 확인됨.
In this thesis a seed fixes the initial weights together with the
random draws of training traces and regularization noise, and the 30
seeds of a group differ only in the seed
(Section~\ref{sec:method-baseline}).

Write the training procedure as $\theta_s=\mathcal{A}(D,s)$ from a
fixed dataset $D$ and seed $s$ to weights. An evaluation score of
$\theta_s$, such as its NLL $L_s$ on one condition, is then a random
variable over seeds, estimated from $S$ seeds by
\begin{equation}
  \hat{\mu}=\frac{1}{S}\sum_{s=1}^{S}L_s,
  \qquad
  \hat{\sigma}=\Bigl(\frac{1}{S-1}\sum_{s=1}^{S}\bigl(L_s-\hat{\mu}\bigr)^2\Bigr)^{1/2} .
  \label{eq:bg-seed-moments}
\end{equation}
$\hat{\mu}$ describes the procedure on average and $\hat{\sigma}$ how
far one trained model, the thing that is deployed, can fall from it.
The standard error $\hat{\sigma}/\sqrt{S}$ shrinks with more seeds;
$\hat{\sigma}$ does not, which is why WILDS reports the standard
deviation across replicates~\cite{koh2021wilds}. Our group score is
$\hat{\mu}$ over 30 seeds and the seed standard deviation is
$\hat{\sigma}$ (Section~\ref{sec:eval-metrics}).
% 원문: Koh et al. Table 1 캡션 p.7 "standard deviation across 3+ replicates"; Bouthillier et al. §2 p.2
%   "realization of a random variable", §5 p.9 "Randomize as many sources of variations as possible". 확인됨.

The spread can be large even when the sources are small: Summers and
Dinneen found that changing a single bit of one initial weight gave
models that differed widely in accuracy after a few epochs, and that
pairs of models disagreed on individual top-1
predictions~\cite{summers2021nondeterminism}. Underspecification names
the case where a pipeline returns many predictors with equivalent
held-out performance in the training domain that behave differently
under deployment conditions~\cite{damour2020underspecification}. The
held-out score cannot tell them apart, so the choice among them is
left to the seed (Section~\ref{sec:related-underspec}).
% 원문: Summers & Dinneen 초록 p.1 "one-bit changes in initial parameters", §3.1 p.3 disagreement;
%   D'Amour et al. 초록 "many predictors with equivalently strong held-out performance in the training domain". 확인됨.

\section{Ensembles and knowledge distillation}
\label{sec:bg-distill}

An ensemble combines $M$ trained members $p_1,\dots,p_M$, usually by the
arithmetic mean of their probabilities, as in deep
ensembles~\cite{lakshminarayanan2016simple}, or by the normalized
geometric mean, the logarithmic opinion pool~\cite{heskes1997selecting}:
\begin{equation}
  \bar{p}(a\mid x)=\frac{1}{M}\sum_{m=1}^{M}p_m(a\mid x),
  \qquad
  \tilde{p}(a\mid x)=\frac{1}{Z(x)}
  \exp\Bigl(\frac{1}{M}\sum_{m=1}^{M}\log p_m(a\mid x)\Bigr) .
  \label{eq:bg-pools}
\end{equation}
$\bar{p}$ minimizes $\sum_m\mathrm{KL}(p_m\,\|\,q)$ and $\tilde{p}$
minimizes $\sum_m\mathrm{KL}(q\,\|\,p_m)$~\cite{hinton2015distilling}.
Our method uses $\bar{p}$ as its target (Equation~\ref{eq:cd-target})
and our post hoc analysis uses $\tilde{p}$
(Equation~\ref{eq:decomposition}). Random initialization and data
shuffling alone gave members that made a good deep
ensemble~\cite{lakshminarayanan2016simple}, so those members differ,
as our seeds do, only in training randomness.
% 원문: Lakshminarayanan et al. (NIPS 2017) §2.4 p.4 "averaging the predicted probabilities", "random
%   initialization ... along with random shuffling ... was sufficient"; Hinton et al. §5.4 pp.6–7 "either the
%   arithmetic or geometric mean, depending on whether we use KL(p, q) or KL(q, p)". 확인됨.

An ensemble of $M$ networks costs $M$ forward passes. Bucilu\v{a} et
al.\ used an ensemble to label a large unlabeled set and trained one
network on those labels~\cite{bucilua2006model}, and Hinton et al.\
generalized this into knowledge distillation, in which the student
matches the teacher's softened distribution $p^{\mathrm{t}}_T$,
usually together with the cross-entropy on the true labels:
\begin{equation}
  \mathcal{L}_{\mathrm{KD}}
  = \lambda_{\mathrm{hard}}\bigl[-\log p_{\theta,1}(y\mid x)\bigr]
  + \lambda_{\mathrm{soft}}\,T^2\,
    \mathrm{KL}\bigl(p^{\mathrm{t}}_T(\cdot\mid x)\,\|\,p_{\theta,T}(\cdot\mid x)\bigr) .
  \label{eq:bg-kd}
\end{equation}
The factor $T^2$ keeps the two terms in proportion as $T$ changes; the
teacher may be an ensemble, and the transfer set of the soft term may
consist entirely of unlabeled data~\cite{hinton2015distilling}.
% 원문: Buciluǎ et al. (KDD 2006) §1 p.535; Hinton et al. §2 pp.2–3 "scale as 1/T^2", "arithmetic or geometric
%   mean ... as the soft targets", "could consist entirely of unlabeled data". 확인됨.
Our consensus distillation is this loss at $T=1$, with the soft term
computed only on unlabeled environment inputs and the labeled term
replaced by our base objective (Equation~\ref{eq:cd-loss}). Distillation
need not compress: born-again networks, students identical to their
teachers but initialized from a different seed, outperformed those
teachers~\cite{furlanello2018born}. A CD student likewise differs from
its teachers only in its seed and loss, but what we ask of it is a
smaller spread across student seeds on conditions absent from the
labeled data (Section~\ref{sec:related-distillation}).
% 원문: Furlanello et al. (ICML 2018) 초록 p.1607 "parameterized identically", "outperform their teachers";
%   Figure 1 캡션 "initialized from a different random seed". 확인됨.

\section{Distribution shift}
\label{sec:bg-shift}

Under distribution shift the training (source) distribution
$P_S(x,y)$ and the test (target) distribution $P_T(x,y)$ differ. Under
covariate shift the input distribution changes while $P(y\mid x)$ stays
fixed~\cite{sugiyama2007covariate}; under label shift the label marginal
changes while $P(x\mid y)$ stays fixed~\cite{lipton2018detecting}:
\begin{equation}
  \text{covariate shift: } P_T(x,y)=P_T(x)\,P_S(y\mid x),
  \qquad
  \text{label shift: } P_T(x,y)=P_T(y)\,P_S(x\mid y) .
  \label{eq:bg-shift-types}
\end{equation}
% 원문: Sugiyama et al. (JMLR 2007) 초록 p.985; Lipton et al. (ICML 2018) 초록 p.3122, §1 두 분해식. 확인됨.
Under covariate shift, weighting each training loss by $P_T(x)/P_S(x)$
restores consistency for a misspecified model, at the price of
instability~\cite{sugiyama2007covariate}. Domain adaptation theory
bounds the target error of a classifier by its source error, a
divergence between the source and target input distributions that can
be estimated from unlabeled samples, and the error of the best joint
hypothesis, which needs target labels~\cite{bendavid2010theory}; this is
why unlabeled target inputs carry information about the shift. WILDS
collects ten datasets with real shifts, such as across hospitals, and
found standard training substantially worse out of distribution than in
distribution~\cite{koh2021wilds}.
% 원문: Sugiyama et al. §2.2 pp.987–988 IWERM "is consistent", "can be rather unstable"; Ben-David et al.
%   (2010) Theorem 2 p.158, p.152 "estimated from finite, unlabeled samples"; Koh et al. 초록 p.5637. 확인됨.

In this thesis the unseen conditions are executions of the same source
procedure with a benign worker running alongside, a worker that no
labeled training run contains (Section~\ref{sec:conditions}); we do not
assume a pure covariate or label shift. The CD student reads unlabeled
inputs from these conditions during training, so our setting is closer
to adaptation with unlabeled target inputs than to domain
generalization.
